% Section 1 -- Introduction & Motivation (P2)
% Numbers in the findings paragraph come from results/E1_depth.csv and results/baselines_timing.csv.

\section{Introduction}
\label{sec:intro}

High-performance computing (HPC) clusters run many jobs at once, and a scheduler decides which job runs on which compute node.
A poor placement leaves one node overloaded while another sits idle, so every job on the busy node finishes late.
Placement also has to respect interactions between jobs: two memory-bound jobs placed on the same node compete for memory bandwidth and slow each other down, so schedulers let users declare \emph{anti-affinity} rules that keep such jobs apart.
Finding the best placement is a combinatorial problem.
Even the simplest version, splitting $n$ jobs between two nodes so that both finish at the same time, has $2^n$ possible placements and is NP-hard in general.

Quantum computers have been proposed as a new tool for problems of this kind.
The Quantum Approximate Optimization Algorithm (QAOA)~\cite{farhi2014} is one of the leading candidates for near-term quantum devices~\cite{preskill2018}.
It works in three steps.
The problem is first rewritten as an energy function, an Ising Hamiltonian, whose lowest-energy state is the best solution~\cite{lucas2014}.
A short, parameterised quantum circuit then prepares a quantum state that favours low-energy solutions.
Finally, a classical optimiser tunes the circuit's parameters.
Whether QAOA can ever beat the best classical methods is still an open question.

Large enough quantum hardware is not yet available to us, so we study QAOA by simulating it on classical computers.
Simulation is itself an HPC problem: storing the state of $n$ qubits takes $2^n$ complex numbers, so each extra qubit doubles both the memory and the work.
We use graphics processing units (GPUs) through NVIDIA's CUDA-Q platform~\cite{cudaq} to reach problem sizes that would be slow or impossible on a CPU.
HPC therefore appears twice in this project: as the application we schedule for, and as the tool that makes the study possible.

\subsection{Research questions}

We study a simplified version of the scheduling problem: place $n$ jobs on two nodes so that the workloads are as balanced as possible, while keeping conflicting jobs on different nodes.
We ask four questions.

\begin{enumerate}
  \item[\textbf{RQ1}] \textbf{Solution quality.} How close does QAOA get to the optimal schedule, and does it improve as the circuit gets deeper?
  \item[\textbf{RQ2}] \textbf{Classical comparison.} How does QAOA compare with classical methods such as simulated annealing and standard partitioning heuristics?
  \item[\textbf{RQ3}] \textbf{Noise robustness.} How much do realistic hardware errors degrade QAOA's performance?
  \item[\textbf{RQ4}] \textbf{HPC scalability.} How much faster are GPUs than CPUs for simulating these circuits, and how large a problem can we reach?
\end{enumerate}
% TODO(team): RQ3 needs E6 and RQ4 needs E4. If either experiment is dropped, move that question to Limitations.

\subsection{Scope}

This project does not aim to show that quantum computers beat classical ones.
At the problem sizes that can be simulated, classical methods win easily, and we show this directly.
Instead, we aim to provide three things:
\begin{enumerate}
  \item a correct formulation of a real-world scheduling problem for a quantum algorithm;
  \item a solid, reproducible implementation, in which every number and figure is produced by a script from a fixed random seed; and
  \item an honest analysis of where QAOA works, where it struggles, and why.
\end{enumerate}

\subsection{Summary of findings}
% TODO(team): replace with the 4-6 key findings agreed at the W11 Monday meeting.

QAOA's solutions improve steadily with circuit depth.
On the median instance, QAOA's average output reaches 83--87\,\% of the way from the worst to the best schedule with one layer, and 94--97\,\% with six layers, for 8 to 20 jobs.
It also finds an exact optimum far more often than random guessing.
With six layers, 14\,\% of its samples are optimal at $n = 10$, about 34 times the rate of random guessing.
At $n = 20$ only 1.4\,\% of samples are optimal, but random guessing is about 7{,}600 times less likely to succeed.
The advantage over random guessing grows with problem size, while the absolute success rate falls.
Classical methods are much stronger at these sizes: an exact solver (CP-SAT) finds and proves the optimum for 28 jobs in about 0.01 seconds.

\subsection{Outline}

Section~\ref{sec:background} introduces Ising models, QAOA and GPU simulation.
Section~\ref{sec:formulation} turns the scheduling problem into a Hamiltonian.
Sections~\ref{sec:method} and~\ref{sec:implementation} describe the QAOA method and its CUDA-Q implementation, and Section~\ref{sec:setup} describes the experiments.
Section~\ref{sec:results} presents the results, Section~\ref{sec:discussion} discusses them and their limitations, and Section~\ref{sec:conclusion} concludes.
